% problem_id: S001-theorem-quillen-spectral-sequence
% source_id: S001 (Stacks Project)
% source_file: cotangent.tex
% statement_locator: cotangent.tex lines 2162--2177
% proof_locator: cotangent.tex lines 2179--2310
% env: theorem
% id_note: label 在本来源内唯一
% status: complete (statement + author proof verbatim + reason + locators)
%
\begin{theorem}[Quillen spectral sequence]
\label{theorem-quillen-spectral-sequence}
Let $A \to B$ be a surjective ring map. Consider the sheaf
$\Omega = \Omega_{\mathcal{O}/A} \otimes_\mathcal{O} \underline{B}$ of
$\underline{B}$-modules on $\mathcal{C}_{B/A}$, see
Section \ref{section-compute-L-pi-shriek}.
Then there is a spectral sequence with $E_1$-page
$$
E_1^{p, q} =
H_{- p - q}(\mathcal{C}_{B/A}, \text{Sym}^p_{\underline{B}}(\Omega))
\Rightarrow \text{Tor}^A_{- p - q}(B, B)
$$
with $d_r$ of bidegree $(r, -r + 1)$.
Moreover, $H_i(\mathcal{C}_{B/A}, \text{Sym}^k_{\underline{B}}(\Omega)) = 0$
for $i < k$.
\end{theorem}

\begin{proof}
Let $I \subset A$ be the kernel of $A \to B$. Let
$\mathcal{J} \subset \mathcal{O}$
be the kernel of $\mathcal{O} \to \underline{B}$. Then
$I\mathcal{O} \subset \mathcal{J}$. Set
$\mathcal{K} = \mathcal{J}/I\mathcal{O}$ and
$\overline{\mathcal{O}} = \mathcal{O}/I\mathcal{O}$.

\medskip\noindent
For every object $U = (P \to B)$ of $\mathcal{C}_{B/A}$
we can choose an isomorphism $P = A[E]$ such that the map
$P \to B$ maps each $e \in E$ to zero. Then
$J = \mathcal{J}(U) \subset P = \mathcal{O}(U)$
is equal to $J = IP + (e; e \in E)$. Moreover
$\overline{\mathcal{O}}(U) = B[E]$ and $K = \mathcal{K}(U) = (e; e \in E)$
is the ideal generated by the variables in the polynomial ring $B[E]$.
In particular it is clear that
$$
K/K^2 \xrightarrow{\text{d}} \Omega_{P/A} \otimes_P B
$$
is a bijection. In other words, $\Omega = \mathcal{K}/\mathcal{K}^2$
and $\text{Sym}_B^k(\Omega) = \mathcal{K}^k/\mathcal{K}^{k + 1}$.
Note that $\pi_!(\Omega) = \Omega_{B/A} = 0$ (Lemma \ref{lemma-identify-H0})
as $A \to B$ is surjective
(Algebra, Lemma \ref{algebra-lemma-trivial-differential-surjective}).
By Lemma \ref{lemma-vanishing-symmetric-powers} we conclude that
$$
H_i(\mathcal{C}_{B/A}, \mathcal{K}^k/\mathcal{K}^{k + 1}) =
H_i(\mathcal{C}_{B/A}, \text{Sym}^k_{\underline{B}}(\Omega)) = 0
$$
for $i < k$. This proves the final statement of the theorem.

\medskip\noindent
The approach to the theorem is to note that
$$
B \otimes_A^\mathbf{L} B = L\pi_!(\mathcal{O}) \otimes_A^\mathbf{L} B =
L\pi_!(\mathcal{O} \otimes_{\underline{A}}^\mathbf{L} \underline{B}) =
L\pi_!(\overline{\mathcal{O}})
$$
The first equality by Lemma \ref{lemma-apply-O-B-comparison},
the second equality by
Cohomology on Sites, Lemma \ref{sites-cohomology-lemma-change-of-rings}, and
the third equality as $\mathcal{O}$ is flat over $\underline{A}$.
The sheaf $\overline{\mathcal{O}}$ has a filtration
$$
\ldots \subset
\mathcal{K}^3 \subset
\mathcal{K}^2 \subset
\mathcal{K} \subset
\overline{\mathcal{O}}
$$
This induces a filtration $F$ on a complex $C$ representing
$L\pi_!(\overline{\mathcal{O}})$ with $F^pC$ representing
$L\pi_!(\mathcal{K}^p)$ (construction of $C$ and $F$ omitted).
Consider the spectral sequence of
Homology, Section \ref{homology-section-filtered-complex}
associated to $(C, F)$. It has $E_1$-page
$$
E_1^{p, q} = H_{- p - q}(\mathcal{C}_{B/A}, \mathcal{K}^p/\mathcal{K}^{p + 1})
\quad\Rightarrow\quad
H_{- p - q}(\mathcal{C}_{B/A}, \overline{\mathcal{O}}) = 
\text{Tor}_{- p - q}^A(B, B)
$$
and differentials $E_r^{p, q} \to E_r^{p + r, q - r + 1}$. To show convergence
we will show that for every $k$ there exists a $c$ such that
$H_i(\mathcal{C}_{B/A}, \mathcal{K}^n) = 0$
for $i < k$ and $n > c$\footnote{A posteriori
the ``correct'' vanishing $H_i(\mathcal{C}_{B/A}, \mathcal{K}^n) = 0$ for
$i < n$ can be concluded.}.

\medskip\noindent
Given $k \geq 0$ set $c = k^2$. We claim that
$$
H_i(\mathcal{C}_{B/A}, \mathcal{K}^{n + c}) \to
H_i(\mathcal{C}_{B/A}, \mathcal{K}^n)
$$
is zero for $i < k$ and all $n \geq 0$. Note that
$\mathcal{K}^n/\mathcal{K}^{n + c}$ has a finite filtration whose successive
quotients $\mathcal{K}^m/\mathcal{K}^{m + 1}$, $n \leq m < n + c$
have $H_i(\mathcal{C}_{B/A}, \mathcal{K}^m/\mathcal{K}^{m + 1}) = 0$
for $i < n$ (see above). Hence the claim implies
$H_i(\mathcal{C}_{B/A}, \mathcal{K}^{n + c}) = 0$ for $i < k$ and all
$n \geq k$ which is what we need to show.

\medskip\noindent
Proof of the claim. Recall that for any $\mathcal{O}$-module $\mathcal{F}$
the map $\mathcal{F} \to \mathcal{F} \otimes_\mathcal{O}^\mathbf{L} B$
induces an isomorphism on applying $L\pi_!$, see
Lemma \ref{lemma-O-homology-B-homology}.
Consider the map
$$
\mathcal{K}^{n + k} \otimes_\mathcal{O}^\mathbf{L} B
\longrightarrow
\mathcal{K}^n \otimes_\mathcal{O}^\mathbf{L} B
$$
We claim that this map induces the zero map on cohomology sheaves
in degrees $0, -1, \ldots, - k + 1$. If this second claim holds, then
the $k$-fold composition
$$
\mathcal{K}^{n + c} \otimes_\mathcal{O}^\mathbf{L} B
\longrightarrow
\mathcal{K}^n \otimes_\mathcal{O}^\mathbf{L} B
$$
factors through $\tau_{\leq -k}\mathcal{K}^n \otimes_\mathcal{O}^\mathbf{L} B$
hence induces zero on $H_i(\mathcal{C}_{B/A}, -) = L_i\pi_!( - )$
for $i < k$, see
Derived Categories, Lemma \ref{derived-lemma-trick-vanishing-composition}.
By the remark above this means the same thing is true for
$H_i(\mathcal{C}_{B/A}, \mathcal{K}^{n + c}) \to
H_i(\mathcal{C}_{B/A}, \mathcal{K}^n)$
which proves the (first) claim.

\medskip\noindent
Proof of the second claim. The statement is local, hence we may work
over an object $U = (P \to B)$ as above. We have to show
the maps
$$
\text{Tor}_i^P(B, K^{n + k}) \to \text{Tor}_i^P(B, K^n)
$$
are zero for $i < k$. There is a spectral sequence
$$
\text{Tor}_a^P(P/IP, \text{Tor}_b^{P/IP}(B, K^n))
\Rightarrow
\text{Tor}_{a + b}^P(B, K^n),
$$
see More on Algebra, Example \ref{more-algebra-example-tor-change-rings}.
Thus it suffices to prove the maps
$$
\text{Tor}_i^{P/IP}(B, K^{n + 1}) \to \text{Tor}_i^{P/IP}(B, K^n)
$$
are zero for all $i$. This is Lemma \ref{lemma-map-tors-zero}.
\end{proof}
